\section{A gate as a substitutable contract}
\label{sec:contract}

A DGF is an organization's composition of review checkpoints. It is not a universal sequence:
enterprises can merge, split, or reorder gates while preserving their obligations and information
dependencies. Typical tasks include supplier due diligence, contractual checks, technical
architecture review, security assessment, readiness verification, and final consolidation.
We study three example routes, Buy, Integrate, and Build; the argument applies to the review
contract, not to these particular names.

\begin{definition}[Review contract]
A gate contract specifies permitted evidence and tools, a versioned policy, an authorization
mandate, and the conditions for accepting an output. The output contains a disposition,
findings, required actions, supporting evidence, and an authorization record. Its accepted
form and substantive requirements are fixed before evaluating a proposed replacement.
\end{definition}

In DGF-Bench the dispositions are \texttt{GO}, \texttt{GO\_WITH\_RESERVATIONS}, \texttt{REWORK},
\texttt{SUSPENSION}, and \texttt{NO\_GO}. A refusal can be a correct completed review. Conversely,
an approval is unsuccessful if it omits a mandatory defect or exceeds the reviewer's mandate.
Completing the review is distinct from implementing the remediation it requests. The same
boundary must apply to human and software work when estimating substitution.

\subsection{Three conditions for a replacement}

\textbf{Information.} The permitted observations must contain what the policy needs, or the
contract must allow a request for information or abstention. A model cannot recover a missing
fact merely by generating a plausible answer. Source accessibility matters as much as physical
file existence; the agent may lack permission to read a relevant record.

\textbf{Validity.} The system must produce outputs satisfying the review's quality and service
requirements on its declared operating population. These include evidence support and the
handling of uncertainty, not just a disposition matching a label. A machine-readable schema
can reject malformed records but cannot by itself prove the substantive decision correct.

\textbf{Authority.} A valid mandate must permit the system's action. An agent's proposal and
the mechanism that makes a commitment effective are separate. Deterministic enforcement can
block forbidden commitments even when an agent requests them; a final authorized output does
not establish that the model respected authority unaided. Delegating bounded execution does
not eliminate the organization's accountability or resolve deployment-specific legal requirements.

These conditions describe functional replacement. Net substitution of human work additionally
requires a reduction in all human hours used to deliver that function, as derived next. Removing
a routine reviewer while adding equivalent effort to project teams or support providers moves
the work rather than reducing it.

\subsection{An information obstruction}

Let $I(x)$ denote all observations permitted in an evaluation condition for underlying case $x$,
and $V(x)$ the set of outputs accepted by its fixed contract. Policy and mandate are held fixed.

\begin{proposition}[Indistinguishable cases]
If two admissible cases satisfy $I(x_1)=I(x_2)$ but
$V(x_1)\cap V(x_2)=\varnothing$, no system using only $I$ can guarantee a valid output on both.
\end{proposition}
\begin{proof}
Equal observations imply equal output distributions, including for a deterministic system.
Guaranteed validity on both cases would require one distribution to assign probability one
to both disjoint accepted sets, which is impossible.
\end{proof}

This elementary observation is an operational test for an automation boundary, not a new general
impossibility theorem. Our offline counterexample makes it concrete: two variants have identical
extracted text in all 26 Word documents, but changing the underlying supplier due-diligence
status, recorded in CSV evidence, changes the fixed reference decision from GO to REWORK.
A condition restricted to those Word
texts therefore cannot reproduce both reference decisions. Allowing the CSV, or changing the
contract to accept a justified request for missing information, changes the task. For the fixed
reference test, the accepted disposition sets are disjoint. This obstruction applies to a human
reader as well as a model; it identifies a deficient information interface.

\subsection{Why the DGF is a candidate}

Table~\ref{tab:candidate} turns the claim that DGF is suitable for automation into properties
that can be checked in an organization. These are engineering advantages when present, not
guarantees that every enterprise already has clean rules and accessible evidence.

\begin{table}[!htb]
\centering\small
\caption{Properties that make a governance workflow a candidate for task substitution.}
\label{tab:candidate}
\begin{tabularx}{\linewidth}{@{}P{3.1cm}LL@{}}
\toprule
Property & Potential advantage & Condition that can defeat it\\
\midrule
Recurring dossiers and outputs & Reusable extraction and decision interfaces & Material facts remain tacit or inaccessible\\
Explicit policies and mandates & Executable checks and bounded delegation & Conflicting rules or discretionary commitments\\
Recorded gate handoffs & Shared evidence and traceable dependencies & Work is repeated or lost between teams\\
Repeated review populations & Integration cost spread across cases & Low volume or frequent policy change\\
Observable acceptance criteria & Regression tests and monitored service & Evaluation rewards format instead of valid work\\
\bottomrule
\end{tabularx}
\end{table}
